% Generated from paper/full_draft. Edit the Markdown source or migration script.

\section{Limitations and Future Work}
\label{sec:08-limitations:limitations-and-future-work}

\subsection{Simulated Feedback}
\label{sec:08-limitations:simulated-feedback}

The current experiments use simulated feedback derived from ground truth and controlled noise/trust settings. This validates whether the policy can improve under a feedback signal, but it does not prove the same behavior under real human feedback. Real deployments must handle delayed feedback, biased implicit signals, adversarial or low-quality users, and non-stationary intent.

\subsection{Limited Generation Evaluation}
\label{sec:08-limitations:limited-generation-evaluation}

The main experiments evaluate retrieval and final retrieved context tokens. A 60-query LLM generation smoke test shows no obvious answer-quality degradation from conservative context compaction, but this is not a full end-to-end human evaluation. The supported main claim remains evidence retrieval and retrieved context budget, not generated answer superiority or user satisfaction.

\subsection{Dense Remains Strong}
\label{sec:08-limitations:dense-remains-strong}

Dense-only retrieval remains a strong baseline. IntentWeight should not be claimed as a universal replacement for dense retrieval. The evidence supports a controller that can reduce final context tokens while preserving dense-level $\mathrm{Hit@10}$ in the main LoTTE setting, and that can expose a quality-cost frontier across routes.

\subsection{Evidence Completeness Trade-Off}
\label{sec:08-limitations:evidence-completeness-trade-off}

The main retrieval headline is query-level $\mathrm{Hit@10}$. Final context compaction can preserve whether at least one relevant chunk is retrieved while reducing the fraction of all ground-truth chunks retrieved. For tasks that require complete evidence collection, such as legal review, medical evidence synthesis, or exhaustive compliance analysis, a more conservative context policy or no compaction may be preferable.

\subsection{Geometry Is Diagnostic, Not a Proof}
\label{sec:08-limitations:geometry-is-diagnostic-not-a-proof}

The piecewise relevance-manifold framing is supported by diagnostics such as $\mathrm{NearestClusterHit@3}$, PCA spectrum, and context retention. These diagnostics do not prove a mathematical manifold theorem. They show that local geometry is informative for routing on LoTTE, while dense retrieval remains necessary.

\subsection{KMeans Is an Experimental Arm Design}
\label{sec:08-limitations:kmeans-is-an-experimental-arm-design}

KMeans/MiniBatchKMeans is used because LinUCB requires a fixed arm space and the experiments need reproducible, scalable arms. This is not a claim that KMeans is the best clustering method for all RAG systems. HDBSCAN or graph-based clusters may perform better in some deployments, but dynamic arm counts complicate the current LinUCB setup.

\subsection{Limited Encoder and Domain Coverage}
\label{sec:08-limitations:limited-encoder-and-domain-coverage}

The main dense baseline uses \texttt{sentence-transformers/all-MiniLM-L6-v2}. The paper adds a CPU-friendly encoder robustness check with \texttt{sentence-transformers/multi-qa-MiniLM-L6-cos-v1}, but the paper should not generalize the result to stronger domain-specific encoders, rerankers, or late-interaction models without additional experiments.

LoTTE technology/search is the main positive large-scale domain. Additional LoTTE domains or other vertical corpora would strengthen external validity.

\subsection{Seed Count and 400k Variance}
\label{sec:08-limitations:seed-count-and-400k-variance}

The stability analysis reports three-seed confidence intervals across LoTTE 100k-638k, and an additional robustness check extends LoTTE 100k to five seeds. These are useful engineering stability diagnostics, but they should not be over-framed as strong statistical significance proof. The LoTTE 400k token-saving interval is notably wider than the other scales and should be interpreted as seed-level variance in route confidence and context-budget control.

\subsection{Future Work}
\label{sec:08-limitations:future-work}

Future work should evaluate:

\begin{itemize}
\item real user feedback with trust scoring and delayed-feedback handling;
\item larger end-to-end LLM answer-quality and citation-faithfulness studies;
\item stronger dense encoders, rerankers, and late-interaction retrieval models;
\item graph or density-based dynamic clustering under bandit-compatible arm management;
\item larger seed counts and additional vertical-domain corpora;
\item production policies that lower dense usage only after route confidence is demonstrably stable.
\end{itemize}
